% Options for packages loaded elsewhere
\PassOptionsToPackage{unicode}{hyperref}
\PassOptionsToPackage{hyphens}{url}
\documentclass[
]{article}
\usepackage{xcolor}
\usepackage{amsmath,amssymb}
\setcounter{secnumdepth}{-\maxdimen} % remove section numbering
\usepackage{iftex}
\ifPDFTeX
  \usepackage[T1]{fontenc}
  \usepackage[utf8]{inputenc}
  \usepackage{textcomp} % provide euro and other symbols
\else % if luatex or xetex
  \usepackage{unicode-math} % this also loads fontspec
  \defaultfontfeatures{Scale=MatchLowercase}
  \defaultfontfeatures[\rmfamily]{Ligatures=TeX,Scale=1}
\fi
\usepackage{lmodern}
\ifPDFTeX\else
  % xetex/luatex font selection
\fi
% Use upquote if available, for straight quotes in verbatim environments
\IfFileExists{upquote.sty}{\usepackage{upquote}}{}
\IfFileExists{microtype.sty}{% use microtype if available
  \usepackage[]{microtype}
  \UseMicrotypeSet[protrusion]{basicmath} % disable protrusion for tt fonts
}{}
\makeatletter
\@ifundefined{KOMAClassName}{% if non-KOMA class
  \IfFileExists{parskip.sty}{%
    \usepackage{parskip}
  }{% else
    \setlength{\parindent}{0pt}
    \setlength{\parskip}{6pt plus 2pt minus 1pt}}
}{% if KOMA class
  \KOMAoptions{parskip=half}}
\makeatother
\usepackage{color}
\usepackage{fancyvrb}
\newcommand{\VerbBar}{|}
\newcommand{\VERB}{\Verb[commandchars=\\\{\}]}
\DefineVerbatimEnvironment{Highlighting}{Verbatim}{commandchars=\\\{\}}
% Add ',fontsize=\small' for more characters per line
\newenvironment{Shaded}{}{}
\newcommand{\AlertTok}[1]{\textcolor[rgb]{1.00,0.00,0.00}{\textbf{#1}}}
\newcommand{\AnnotationTok}[1]{\textcolor[rgb]{0.38,0.63,0.69}{\textbf{\textit{#1}}}}
\newcommand{\AttributeTok}[1]{\textcolor[rgb]{0.49,0.56,0.16}{#1}}
\newcommand{\BaseNTok}[1]{\textcolor[rgb]{0.25,0.63,0.44}{#1}}
\newcommand{\BuiltInTok}[1]{\textcolor[rgb]{0.00,0.50,0.00}{#1}}
\newcommand{\CharTok}[1]{\textcolor[rgb]{0.25,0.44,0.63}{#1}}
\newcommand{\CommentTok}[1]{\textcolor[rgb]{0.38,0.63,0.69}{\textit{#1}}}
\newcommand{\CommentVarTok}[1]{\textcolor[rgb]{0.38,0.63,0.69}{\textbf{\textit{#1}}}}
\newcommand{\ConstantTok}[1]{\textcolor[rgb]{0.53,0.00,0.00}{#1}}
\newcommand{\ControlFlowTok}[1]{\textcolor[rgb]{0.00,0.44,0.13}{\textbf{#1}}}
\newcommand{\DataTypeTok}[1]{\textcolor[rgb]{0.56,0.13,0.00}{#1}}
\newcommand{\DecValTok}[1]{\textcolor[rgb]{0.25,0.63,0.44}{#1}}
\newcommand{\DocumentationTok}[1]{\textcolor[rgb]{0.73,0.13,0.13}{\textit{#1}}}
\newcommand{\ErrorTok}[1]{\textcolor[rgb]{1.00,0.00,0.00}{\textbf{#1}}}
\newcommand{\ExtensionTok}[1]{#1}
\newcommand{\FloatTok}[1]{\textcolor[rgb]{0.25,0.63,0.44}{#1}}
\newcommand{\FunctionTok}[1]{\textcolor[rgb]{0.02,0.16,0.49}{#1}}
\newcommand{\ImportTok}[1]{\textcolor[rgb]{0.00,0.50,0.00}{\textbf{#1}}}
\newcommand{\InformationTok}[1]{\textcolor[rgb]{0.38,0.63,0.69}{\textbf{\textit{#1}}}}
\newcommand{\KeywordTok}[1]{\textcolor[rgb]{0.00,0.44,0.13}{\textbf{#1}}}
\newcommand{\NormalTok}[1]{#1}
\newcommand{\OperatorTok}[1]{\textcolor[rgb]{0.40,0.40,0.40}{#1}}
\newcommand{\OtherTok}[1]{\textcolor[rgb]{0.00,0.44,0.13}{#1}}
\newcommand{\PreprocessorTok}[1]{\textcolor[rgb]{0.74,0.48,0.00}{#1}}
\newcommand{\RegionMarkerTok}[1]{#1}
\newcommand{\SpecialCharTok}[1]{\textcolor[rgb]{0.25,0.44,0.63}{#1}}
\newcommand{\SpecialStringTok}[1]{\textcolor[rgb]{0.73,0.40,0.53}{#1}}
\newcommand{\StringTok}[1]{\textcolor[rgb]{0.25,0.44,0.63}{#1}}
\newcommand{\VariableTok}[1]{\textcolor[rgb]{0.10,0.09,0.49}{#1}}
\newcommand{\VerbatimStringTok}[1]{\textcolor[rgb]{0.25,0.44,0.63}{#1}}
\newcommand{\WarningTok}[1]{\textcolor[rgb]{0.38,0.63,0.69}{\textbf{\textit{#1}}}}
\setlength{\emergencystretch}{3em} % prevent overfull lines
\providecommand{\tightlist}{%
  \setlength{\itemsep}{0pt}\setlength{\parskip}{0pt}}
\usepackage{bookmark}
\IfFileExists{xurl.sty}{\usepackage{xurl}}{} % add URL line breaks if available
\urlstyle{same}
\hypersetup{
  pdftitle={Formalization workflow for U-AUB-02},
  pdfauthor={Kieran McShane},
  hidelinks,
  pdfcreator={LaTeX via pandoc}}

\title{Formalization workflow for U-AUB-02}
\author{Kieran McShane}
\date{2026-06-07}

\begin{document}
\maketitle

\section{Formalization workflow for
U-AUB-02}\label{formalization-workflow-for-u-aub-02}

\subsection{Source of truth}\label{source-of-truth}

The Lean source remains the authority. Article prose, blueprint
statements, and ticket reports must be generated from checked
declarations, board entries, and axiom audits.

Current main module:

\begin{Shaded}
\begin{Highlighting}[]
\NormalTok{PptFactorization.AppendixBAubrunGraduate}
\end{Highlighting}
\end{Shaded}

Global status:

\begin{Shaded}
\begin{Highlighting}[]
\NormalTok{U{-}AUB{-}02 is open.  The active endpoint is a checked conditional reduction,}
\NormalTok{not a no{-}input Aubrun count theorem.}
\end{Highlighting}
\end{Shaded}

Active theorem-facing endpoint:

\begin{Shaded}
\begin{Highlighting}[]
\NormalTok{AppendixB.AubrunBiDefectCountBound\_of\_bianeAdjacentLinkHalves\_left\_gap\_three\_start\_impossible\_frontiers\_nestedStrictReturnIntervalOpenNonwrapOffCycleSmallerInternalOnly\_and\_chapuy}
\end{Highlighting}
\end{Shaded}

This is the current reduction layer to track while U-AUB-02 remains
open. It is sharper than the previous primitive-child frontier because
the five strict smaller child primitive leaves are now supplied
internally by the parent primitive residual input. The remaining smaller
nonwrapping child leaf is the genuine internal-subinterval collision
branch.

Previous reduction layer:

\begin{Shaded}
\begin{Highlighting}[]
\NormalTok{AppendixB.AubrunBiDefectCountBound\_of\_bianeAdjacentLinkHalves\_left\_gap\_three\_start\_impossible\_frontiers\_nestedStrictReturnIntervalOpenNonwrapOffCycleSmallerPrimitiveBranchCases\_and\_chapuy}
\end{Highlighting}
\end{Shaded}

Delta: the previous layer still exposed internal \texttt{S}, start
\texttt{S}, end \texttt{S}, fixed predecessor, adjacent \texttt{T}, and
genuine shorter internal-subinterval as separate smaller child leaves.
The active layer uses
\texttt{leftTightStrictReturnIntervalPrimitiveResidual\_of\_strictChildPrimitive}
to absorb the five primitive alternatives into the parent primitive
leaf.

Current theorem-facing surface:

Every item in the following list is still an input to the active
endpoint, not a discharged subgoal. All finite gap-three branches are
omitted because they are now supplied inside the endpoint.

\begin{Shaded}
\begin{Highlighting}[]
\NormalTok{large{-}gap left Biane link, only for 2 \textless{} c{-}a and c{-}a != 3}
\NormalTok{large{-}gap right Biane link, outside the checked d{-}b=3, c=b+2 slice}
\NormalTok{five unchanged side{-}specific residual branch collision leaves}
\NormalTok{one smaller nonwrapping child internal{-}subinterval collision leaf}
\NormalTok{Chapuy/trisection recurrence}
\end{Highlighting}
\end{Shaded}

Presentation rule: do not cite the active endpoint as ``the Aubrun
theorem'' without the qualifier ``conditional reduction from the
remaining theorem-facing leaves''. The name contains many checked
components, including \texttt{chapuy}, but Chapuy is still a hypothesis
of the endpoint, not an internal proof.

The left-start slice \texttt{c-a=3,\ b=a+1} no longer appears on this
surface. Its four former branches are discharged internally: one from
the right frontier, two by forcing that right-frontier contradiction,
and the residual \texttt{not\ pi(a,a+2)\ and\ S(a,a+2)} branch by tree
uniqueness for the copied length-two contour path.

\subsection{Scientific toolchain}\label{scientific-toolchain}

The local proof workflow now has a single status command:

\begin{Shaded}
\begin{Highlighting}[]
\NormalTok{python3 tools/scientific\_tool\_status.py}
\end{Highlighting}
\end{Shaded}

It checks Python, Lake, Graphviz, Pandoc, latexdiff, leanblueprint, the
leanblueprint plasTeX runtime, Grok, optional narration availability,
the reference-library bridge, and the blueprint source files. The
command reports readiness only; it does not print secrets.

Finite permutation exploration is handled by:

\begin{Shaded}
\begin{Highlighting}[]
\NormalTok{python3 scripts/u\_aub\_02\_perm\_explorer.py {-}{-}max{-}n 7}
\NormalTok{python3 scripts/aubrun\_branch\_sanity.py {-}{-}max{-}n 9 {-}{-}all}
\end{Highlighting}
\end{Shaded}

These scripts are conjecture-finding and regression tooling only. Their
outputs cannot close a Lean theorem-facing leaf.

The proof-frontier picture is generated by:

\begin{Shaded}
\begin{Highlighting}[]
\NormalTok{python3 tools/u\_aub\_02\_frontier\_graph.py}
\end{Highlighting}
\end{Shaded}

It writes \texttt{build/u\_aub\_02\_frontier.\{dot,svg,png\}} and
deliberately marks the four live leaves as open.

The maintained artifact build is:

\begin{Shaded}
\begin{Highlighting}[]
\NormalTok{tools/build\_formalization\_artifacts.sh}
\end{Highlighting}
\end{Shaded}

This rebuilds the status graph, blueprint PDF/web output, blueprint
declaration checks, and the Pandoc/LaTeX exposition PDFs. The script
pins the leanblueprint plasTeX runtime; using the global
\texttt{plastex} binary can silently omit the leanblueprint plugins.

Zotero is available through:

\begin{Shaded}
\begin{Highlighting}[]
\NormalTok{tools/zotero\_project\_refs.sh status}
\NormalTok{tools/zotero\_project\_refs.sh search "Aubrun"}
\NormalTok{tools/zotero\_project\_refs.sh export{-}all}
\end{Highlighting}
\end{Shaded}

The reference bridge is optional. Bibliography writes or imports should
still be treated as explicit project changes.

Sage/GAP is optional. Use it for larger finite-combinatorics experiments
if it is available:

\begin{Shaded}
\begin{Highlighting}[]
\NormalTok{brew install {-}{-}cask sage}
\end{Highlighting}
\end{Shaded}

Checked normalization machinery for the current gap-three surface:

\begin{Shaded}
\begin{Highlighting}[]
\NormalTok{AppendixB.AubrunBiDefectCountBound\_of\_bianeAdjacentLinkHalves\_left\_gap\_three\_start\_impossible\_frontiers\_nestedReturnContourIntervalCommonCycle\_and\_chapuy}
\NormalTok{AppendixB.AubrunBiDefectCountBound\_of\_bianeAdjacentLinkHalves\_left\_gap\_three\_start\_impossible\_frontiers\_skipMinimalFrontierCommonCycle\_and\_chapuy}
\NormalTok{AppendixB.leftTight\_bianeAdjacentLink\_left\_of\_gap\_three\_start\_adjacent\_reduced\_fixed\_a2\_no\_pi\_a\_a2}
\NormalTok{AppendixB.leftTight\_bianeAdjacentLink\_left\_of\_gap\_three\_start\_adjacent\_target\_impossible}
\NormalTok{AppendixB.leftTight\_bianeAdjacentLink\_left\_of\_gap\_three\_start\_adjacent\_branch\_to\_target\_iff\_false}
\NormalTok{AppendixB.leftTight\_bianeAdjacentLink\_left\_of\_gap\_three\_start\_S\_b\_c\_false\_of\_global\_right\_frontier}
\NormalTok{AppendixB.leftTight\_bianeAdjacentLink\_left\_of\_gap\_three\_start\_fixed\_a2\_false\_of\_global\_right\_frontier}
\NormalTok{AppendixB.leftTight\_bianeAdjacentLink\_left\_of\_gap\_three\_start\_T\_b\_a2\_false\_of\_global\_right\_frontier}
\NormalTok{AppendixB.leftTight\_bianeAdjacentLink\_left\_of\_gap\_three\_start\_S\_a\_a2\_no\_pi\_false}
\NormalTok{AppendixB.leftTight\_bianeAdjacentLink\_right\_of\_gap\_three\_end\_adjacent\_of\_pi\_b\_c}
\NormalTok{AppendixB.leftTight\_bianeAdjacentLink\_right\_of\_gap\_three\_end\_adjacent\_of\_S\_b\_c}
\NormalTok{AppendixB.leftTight\_bianeAdjacentLink\_right\_of\_gap\_three\_end\_adjacent\_T\_b1\_c\_to\_S\_b1\_d\_of\_right\_rest}
\NormalTok{AppendixB.leftTight\_bianeAdjacentLink\_right\_of\_gap\_three\_end\_adjacent\_fixed\_b1\_to\_S\_b1\_d}
\NormalTok{AppendixB.leftTight\_right\_gap\_three\_T\_b1\_c\_of\_S\_b1\_d\_of\_not\_S\_c\_d}
\NormalTok{AppendixB.leftTight\_bianeAdjacentLink\_right\_of\_gap\_three\_end\_adjacent\_of\_S\_b1\_d\_of\_left\_shift}
\NormalTok{AppendixB.leftTight\_bianeAdjacentLink\_right\_of\_gap\_three\_end\_adjacent\_reduced\_fixed\_b1\_no\_pi\_b\_b1\_no\_pi\_b\_c}
\NormalTok{AppendixB.leftTight\_bianeAdjacentLink\_left\_of\_gap\_three\_end\_adjacent\_reduced}
\NormalTok{AppendixB.leftTight\_bianeAdjacentLink\_left\_of\_gap\_three\_end\_adjacent\_of\_pi\_a\_a1}
\NormalTok{AppendixB.leftTight\_bianeAdjacentLink\_left\_of\_gap\_three\_end\_adjacent\_reduced\_no\_pi\_a\_a1}
\NormalTok{AppendixB.leftTight\_bianeAdjacentLink\_left\_of\_gap\_three\_end\_adjacent\_of\_pi\_a1\_b}
\NormalTok{AppendixB.leftTight\_bianeAdjacentLink\_left\_of\_gap\_three\_end\_adjacent\_reduced\_no\_pi\_branches}
\NormalTok{AppendixB.leftTight\_bianeAdjacentLink\_left\_of\_gap\_three\_end\_adjacent\_reduced\_fixed\_a1\_no\_pi\_branches}
\end{Highlighting}
\end{Shaded}

These helpers justify branch rewrites inside the active endpoint. Their
checked status is not an unconditional Aubrun count theorem.

The checked right gap-three reducer consumes the branch \texttt{pi(b,c)}
by transitivity with the crossing link \texttt{pi(a,c)}. The
adjacent-repeat branch \texttt{S(b+1,c)} has been normalized to the
fixed-point branch \texttt{pi(b+1)=b+1}, and the special branch
\texttt{S(b,c)} is now closed by copying the contour step
\texttt{c\ -\textgreater{}\ d} across \texttt{S(b)=S(c)} and comparing
it with the direct \texttt{pi(b,d)} path. The branch \texttt{pi(b+1,c)}
is now routed through the shifted crossing \texttt{(a,b,b+1,d)}: it
gives \texttt{pi(a,b+1)}, the large-gap right input gives
\texttt{S(b+1,d)}. The fixed branch \texttt{pi(b+1)=b+1} is routed
through the same branch: fixedness gives \texttt{S(b+1,c)}, and denying
\texttt{S(b+1,d)} creates a length-four incidence path from
\texttt{S(b)} to \texttt{S(d)} that contradicts the direct length-two
path from \texttt{pi(b,d)}. The last branch \texttt{S(b+1,d)} is now
closed: if \texttt{S(c,d)} failed, incidence-path uniqueness forces
\texttt{pi(b+1,c)}; then the shifted left link gives \texttt{S(a,b)},
and the incidence-square obstruction contradicts the branch. The active
endpoint signature therefore contains no \texttt{hRightS\_b\_c} or
\texttt{hRightT\_b1\_c} or \texttt{hRightFixed\_b1} or
\texttt{hRightS\_b1\_d} parameter. The left start-adjacent slice
\texttt{c-a=3,\ b=a+1} is fully discharged at the active endpoint. The
left-start \texttt{S(b,c)} usage in that now-discharged slice is
supplied through the global right-frontier wrapper; this is distinct
from the right gap-three special branch \texttt{S(b,c)} above, which is
now discharged locally. The fixed \texttt{a+2} branch and the
\texttt{pi(b,a+2)} branch both force the left-start \texttt{S(b,c)}
contradiction. The residual \texttt{not\ pi(a,a+2)} and
\texttt{S(a,a+2)} branch is closed by copying the contour step from
\texttt{a+2} to \texttt{c} into a length-two incidence path from
\texttt{S(a)} to \texttt{S(c)}; tree uniqueness compares it with the
direct path from \texttt{pi(a,c)} and forces the forbidden link
\texttt{pi(a,a+2)}. The left gap-three end-adjacent slice
\texttt{c-a=3,\ b=a+2} is reduced to two finite branches. The branch
\texttt{pi(a,a+1)} is absorbed by the shifted gap-two left theorem. The
branch \texttt{pi(a+1,b)} is absorbed by the already visible
start-adjacent branch on the shifted crossing \texttt{(a,a+1,c,d)}. The
adjacent-repeat branch \texttt{S(a+1,b)} is normalized to the
fixed-point branch \texttt{pi(a+1)=a+1}; the other theorem-facing
end-adjacent branch left in that slice is \texttt{S(a+1,c)}.

These branch counts are endpoint-local snapshots, not a progress
percentage. Every finite-branch removal must still restate the non-local
leaves separately: the large-gap Biane inputs, the nested non-direct
return-interval collision theorem, and the Chapuy/trisection recurrence.
A wrapper only counts as real progress when the named endpoint exposes
strictly fewer branch hypotheses after the change.

\subsection{Verification tiers}\label{verification-tiers}

Use three tiers, in this order:

\begin{enumerate}
\def\labelenumi{\arabic{enumi}.}
\tightlist
\item
  micro-check a local theorem or signature with
  \texttt{lake\ env\ lean};
\item
  quick-check the enclosing module with
  \texttt{scripts/lean\_quick\_check.py};
\item
  rebuild the enclosing module with \texttt{lake\ build}.
\end{enumerate}

Only after the module build passes should the ledger, board, article
note, or blueprint be updated.

\subsection{Branch-probe discipline}\label{branch-probe-discipline}

The Appendix B file is large and declaration order matters. New Biane
finite branch lemmas should therefore be developed outside the endpoint
first:

\begin{enumerate}
\def\labelenumi{\arabic{enumi}.}
\tightlist
\item
  test the proposed branch implication in small concrete models;
\item
  prove the helper in a downstream Lean probe importing
  \texttt{PptFactorization.AppendixBAubrunGraduate};
\item
  only then move the checked helper into the earlier local branch
  cluster, so the public endpoint never references a declaration that
  appears later in the file.
\end{enumerate}

The finite sanity command is:

\begin{Shaded}
\begin{Highlighting}[]
\NormalTok{python3 scripts/aubrun\_branch\_sanity.py {-}{-}max{-}n 9}
\end{Highlighting}
\end{Shaded}

This script enumerates small left-tight permutations and searches for
counterexamples to the currently exposed branch implications. Since
U-AUB-322 closed the finite gap-three branch list, it now also reports a
live large-gap profile. Through \texttt{n=9}, the profile checks 6,916
left-tight permutations and finds no crossing hypotheses at all. This is
not proof evidence: it supports the expected direct
noncrossing/large-gap route, but it does not close the structural Biane
theorem. Its purpose is only to stop obviously false local branch
attempts before they enter Lean.

\subsection{Article track}\label{article-track}

Narrate Lean is the preferred first pass:

\begin{Shaded}
\begin{Highlighting}[]
\NormalTok{uvx narratelean narrate PptFactorization/AppendixBAubrunGraduate.lean \textbackslash{}}
\NormalTok{  {-}{-}no{-}definitions {-}{-}no{-}proofs \textbackslash{}}
\NormalTok{  {-}o exposition/narratelean\_appendix\_b\_statements.md}
\end{Highlighting}
\end{Shaded}

Current policy: narration is optional exposition support, not proof
evidence. Until a narration pass produces useful mathematics,
leanblueprint plus hand-maintained Pandoc notes are the reliable path.
The note is written directly from theorem statements, proof-artifact
reports, and blueprint declaration checks. Any future Narrate Lean
output still has to be audited against the same distinction: checked
reduction is not no-input closure, and local branch normalization is not
a theorem-strength leaf closure.

Pandoc is the Markdown-to-LaTeX bridge:

\begin{Shaded}
\begin{Highlighting}[]
\NormalTok{pandoc exposition/u\_aub\_02\_formalization\_note.md \textbackslash{}}
\NormalTok{  {-}o exposition/u\_aub\_02\_formalization\_note.tex}
\end{Highlighting}
\end{Shaded}

\subsection{Blueprint track}\label{blueprint-track}

Leanblueprint has been initialized in \texttt{blueprint/}.

Use these conventions:

\begin{enumerate}
\def\labelenumi{\arabic{enumi}.}
\tightlist
\item
  Put mathematical prose in \texttt{blueprint/src/content.tex}.
\item
  Attach Lean declarations using \texttt{\textbackslash{}lean\{...\}}.
\item
  Mark statements with \texttt{\textbackslash{}leanok} only after the
  declaration builds and the axiom audit is benign.
\item
  Keep open theorem-strength leaves visibly unmarked or marked as open
  in prose.
\item
  For conditional reduction theorems, say explicitly that
  \texttt{\textbackslash{}leanok} verifies the reduction from its
  hypotheses, not the hypotheses themselves.
\item
  Every open leaf should name the missing invariant it is expected to
  prove or preserve, not only the Lean placeholder name.
\end{enumerate}

Useful commands once content is stable:

\begin{Shaded}
\begin{Highlighting}[]
\NormalTok{CFLAGS="{-}I/opt/homebrew/include" LDFLAGS="{-}L/opt/homebrew/lib" \textbackslash{}}
\NormalTok{  uvx leanblueprint pdf}
\end{Highlighting}
\end{Shaded}

Then generate the web declaration cache:

\begin{Shaded}
\begin{Highlighting}[]
\NormalTok{CFLAGS="{-}I/opt/homebrew/include" LDFLAGS="{-}L/opt/homebrew/lib" \textbackslash{}}
\NormalTok{  uvx leanblueprint web}
\end{Highlighting}
\end{Shaded}

Finally check that every \texttt{\textbackslash{}lean\{...\}}
declaration referenced from the blueprint exists:

\begin{Shaded}
\begin{Highlighting}[]
\NormalTok{CFLAGS="{-}I/opt/homebrew/include" LDFLAGS="{-}L/opt/homebrew/lib" \textbackslash{}}
\NormalTok{  uvx leanblueprint checkdecls}
\end{Highlighting}
\end{Shaded}

In this repo, \texttt{lake\ build\ EnsX2026} may be needed once before
\texttt{checkdecls}, because the Lake project still advertises that
auxiliary library.

Doc generation through \texttt{doc-gen4} is deferred. The first attempt
tried to switch to \texttt{leanprover/lean4:v4.31.0-rc1} and download a
large toolchain. Do not enable that until the project deliberately
chooses a documentation toolchain policy.

\subsection{LeanArchitect track}\label{leanarchitect-track}

LeanArchitect is the right long-term model: keep blueprint metadata
close to Lean declarations, infer dependencies, and export synchronized
project documents. It is not currently installed in this local repo, so
the immediate workflow is:

\begin{enumerate}
\def\labelenumi{\arabic{enumi}.}
\tightlist
\item
  maintain leanblueprint declarations now;
\item
  keep each blueprint theorem linked to exactly one or a small number of
  Lean constants;
\item
  later migrate those links to LeanArchitect annotations if the package
  is added to the project.
\end{enumerate}

\subsection{Adversarial review loop}\label{adversarial-review-loop}

After any significant proof-frontier or exposition change:

\begin{enumerate}
\def\labelenumi{\arabic{enumi}.}
\tightlist
\item
  ask Grok for adversarial feedback on both the article-style exposition
  and the workflow note;
\item
  treat the feedback as advisory, not as proof evidence;
\item
  patch only claims that can be verified against Lean declarations,
  build output, axiom audits, or generated artifacts;
\item
  rerun Pandoc and leanblueprint checks after incorporating accepted
  changes.
\end{enumerate}

Use the project review command:

\begin{Shaded}
\begin{Highlighting}[]
\NormalTok{tools/grok\_review\_u\_aub\_02.sh}
\end{Highlighting}
\end{Shaded}

It asks for adversarial feedback from the current theorem-frontier
context and keeps the retained Markdown review at
\texttt{build/grok\_u\_aub\_02\_review.md} mathematical.

Durable summaries of useful Grok feedback should be tracked rather than
left only under \texttt{build/}. The current summary is the Grok
workflow review note in the exposition folder. Its main guardrail is to
foreground the four open theorem-strength leaves before the long checked
endpoint name, because a checked conditional adapter can otherwise feel
psychologically close to a done ticket.

Before updating this workflow note or the article note, print the
current theorem header of the sharp
\texttt{AubrunBiDefectCountBound\_of\_*} declaration and confirm that
the listed leaves match its parameters: large-gap left outside
\texttt{c-a=3}, large-gap right outside \texttt{d-b=3} with
\texttt{c=b+2}, no finite gap-three branch parameters, the five
unchanged side-specific residual collision inputs, the smaller
nonwrapping child internal-subinterval input, and \texttt{hstep}.

The recurring review questions are:

\begin{enumerate}
\def\labelenumi{\arabic{enumi}.}
\tightlist
\item
  Does the exposition overclaim a conditional reduction as an
  unconditional theorem?
\item
  Are all theorem-facing leaves named in plain mathematical terms?
\item
  Does every \texttt{\textbackslash{}leanok} marker correspond to a
  declaration that builds and has a benign axiom audit?
\item
  Are the remaining invariants structural, or is the text merely listing
  cases without a proof strategy?
\item
  Does the endpoint-local branch ledger strictly shrink, while the hard
  collision and Chapuy leaves remain visible when they are still open?
\end{enumerate}

\subsection{Exit criteria for
U-AUB-02}\label{exit-criteria-for-u-aub-02}

Do not mark U-AUB-02 done until the public Aubrun expectation pipeline
no longer exposes the corrected count theorem as an input. At minimum,
Lean must check:

\begin{enumerate}
\def\labelenumi{\arabic{enumi}.}
\tightlist
\item
  large-gap left Biane adjacent link outside \texttt{c-a=3};
\item
  large-gap right Biane adjacent link;
\item
  nested non-direct return-interval collision theorem;
\item
  Chapuy/trisection recurrence;
\item
  the final bidefect count endpoint;
\item
  the public expectation pipeline consuming that endpoint.
\end{enumerate}

Each public endpoint must have an axiom audit no worse than:

\begin{Shaded}
\begin{Highlighting}[]
\NormalTok{[propext, Classical.choice, Quot.sound]}
\end{Highlighting}
\end{Shaded}


\end{document}
